% Generated draft: requires mathematical and source review.
\documentclass[11pt]{article}
\usepackage[margin=25mm]{geometry}
\usepackage{fontspec}
\setmainfont{Latin Modern Roman}
\usepackage{amsmath,amssymb,mathtools,unicode-math}
\setmathfont{Latin Modern Math}
\usepackage{enumitem,needspace,xurl,hyperref}
\setlength{\parindent}{0pt}
\setlength{\parskip}{5pt}
\setlength{\emergencystretch}{3em}
\newcommand{\N}{\mathbb{N}}
\newcommand{\Z}{\mathbb{Z}}
\newcommand{\Q}{\mathbb{Q}}
\newcommand{\R}{\mathbb{R}}
\newcommand{\C}{\mathbb{C}}
\newcommand{\F}{\mathbb{F}}
\newcommand{\A}{\mathbb{A}}
\newcommand{\PP}{\mathbb{P}}
\newcommand{\E}{\mathbb{E}}
\newcommand{\eps}{\varepsilon}
\newcommand{\dd}{\,\mathrm{d}}
\begin{document}
\section*{UnsolvedMath ID: 30000113}
\textbf{Source problem code:} catalogue entry
\subsection{Definitions and mathematical statement}
\textbf{Research problem.} Find effective upper bounds for the lattice-point enumerator\[G(K)=\#(K\cap\Z^n)\] in terms of functionals of a convex body \(K\). In the usual convention, a convex body in \(\R^n\) is a compact convex set with nonempty interior, so \(G(K)\) is finite; points on the boundary are included. The source does not explicitly say whether lower-dimensional compact convex sets are also to be covered.\par More precisely, after choosing the permitted geometric or lattice-relative functionals \(F_1,\ldots,F_r\), seek a stated, effectively evaluable bound \(B_n\) such that \(G(K)\le B_n(F_1(K),\ldots,F_r(K))\) for every body in the intended class, with all constants and dimension dependence explicit. The allowable functionals, desired sharpness or asymptotic regime, dependence on lattice position, and any symmetry or integrality hypotheses are not supplied. Consequently this is an upper-bound research program, not one specified numerical inequality. No restriction to origin-symmetric bodies, lattice polytopes, or integral dilations is imposed merely because the report concerns Ehrhart quasipolynomials.\par The supplied literature triage dated 2026-08-21 labels the problem open. It reports subsequent bounds for lattice-point counts and related convex-body functionals, while describing the requested good general upper bound as unresolved. Since \emph{good} and the functionals are not fixed here, that dated assessment is retained as a description of the intended research direction rather than an independently verified claim that no effective upper bounds exist.
\subsection{Short English statement}
Bound the number of integer points, including boundary points, in a convex body using specified geometric or lattice-related data. The excerpt leaves the functionals, body class, and target sharpness open, so it states a research direction rather than a single conjectured inequality.
\subsection{Sources}
\begin{itemize}\item Local UnsolvedMath registry, permanent ID 30000113, code catalogue entry; \emph{Convex-Body Bounds for Lattice-Point Enumerators}, Oberwolfach Reports Open Problems.\item \emph{Ehrhart Quasipolynomials: Algebra, Combinatorics, and Geometry (2005)}, Oberwolfach Reports, DOI 10.4171/owr/2004/39, pp. 2071--2102. \url{https://doi.org/10.4171/owr/2004/39}. Proposed year and year tag: 2004; no proposer named.\item Supplied literature label \emph{OWR report}: \url{https://doi.org/10.4171/owr/2004/39}.\item Supplied literature label \emph{Later lattice-point bounds}: \url{https://link.springer.com/article/10.1007/s00454-021-00310-7}.\item Literature triage checked 2026-08-21, status \texttt{open}; record updated 2026-09-16. The \texttt{corrected\_verified} field reports checking the official report; that check and the dated assessment were not independently verified here.\end{itemize}
\end{document}
